\documentclass{article}
\usepackage{geometry}
\usepackage{graphicx} % Required for inserting images
\usepackage{amsmath, amsthm, amssymb}
\usepackage{parskip}
\newgeometry{vmargin={15mm}, hmargin={24mm,34mm}}
\theoremstyle{definition} 
\newtheorem{definition}{Definition}

\newtheorem{theorem}{Theorem}[section]
\newtheorem{lemma}[theorem]{Lemma}
\newtheorem{corollary}{Corollary}[theorem]


\title{MATH210B Homework 2}
\date{February 2024}
\author{Boran Erol}

\begin{document}

\maketitle

\section{Problem 2}

Let's first show that the map taking a set $ X $ to $ \mathbb{Z} \langle X \rangle $
is left adjoint to the forgetful functor taking a (possibly noncommutative ring $ R $)
to its underlying set.

Let $ \phi : X \to S $ be a set map. Define $ \phi^{*} : \mathbb{Z} \langle X \rangle \to S $
by taking $ x \in X $ to $ \phi(x) $. Let's show that this is a ring homomorphism.

We'll use the following lemma:

\begin{lemma}
    Assume $ A $ is a ring that is also a free Abelian group with basis $ \beta $.
    Assume $ \beta \cdot \beta \subset \beta $.
    Let $ f : A \to B $ be a group homomorphism such that it's a ring homomorphism when restricted
    to $ \beta $. Then, $ f $ is a ring homomorphism. 
\end{lemma}
\begin{proof}
    Let $ x,y \in A $. Then, $ x = \sum_{i = 1}^{n} a_{i}v_{i} $ and $ y = \sum_{i = 1}^{n} b_{i}v_{i} $
    where $ a_{i}, b_{i} \in \mathbb{Z} $ (possibly zero) and $ v_{i} \in \beta $. By distributivity,

    \[ xy = \sum_{i = 1}^{n}\sum_{j = 1} a_{i}b_{j} v_{i}v_{j} \]

    Then,

    \begin{align*}
        f(xy) &= f\Biggl(\sum_{i=1}^{n}\sum_{j=1}^{n} a_i b_j\, v_i v_j\Biggr)\\
              &= \sum_{i=1}^{n}\sum_{j=1}^{n} a_i b_j\, f(v_i v_j)\\
              &= \sum_{i=1}^{n}\sum_{j=1}^{n} a_i b_j\, f(v_i) f(v_j)\\
              &= \sum_{i=1}^{n} f(a_{i}{v_{i}}) \sum_{j=1}^{n} f(a_{j}v_j)\\
              &= f(x)f(y)
    \end{align*}
\end{proof}

Now, notice that $ \phi^{*} $ is an Abelian group since it matches the coproduct in the category
of Abelian groups, and it satisfies the conditions of the lemma. Thus, it's a ring homomorphism.

The fact that these maps are adjoint is now immediate.

Let $ X = R $ as a set. Then, there's a surjective
ring homomorphism from $ \mathbb{Z} \langle X \rangle $ to $ R $
which is given by $ r \mapsto r $. Thus, we have the desired result.

\section{Problem 3}

Let $ A,B $ be (possibly noncommutative) rings.

By Problem 2, we know that $ A \equiv \mathbb{Z}[X] / I $ and  
$ B \equiv \mathbb{Z}[Y] / J $ for some sets $ X,Y $ and ideals $ I,J $
of the ring $ R $.

Let the coproduct of $ A $ and $ B $ be $ \mathbb{Z}[X \sqcup Y ]/ (I + J) $.

We'll show that this is the coproduct by showing that it satisfies the universal property.

Notice that $ \mathbb{Z}[X] / I $ injects into $ \mathbb{Z}[X \sqcup Y]/ (I + J) $
since $ I $ gets mapped to $ I + J $ where we send $ i \mapsto i + 0 $. By a similar argument with $ B $,
we have injective ring homomorphisms into $ \mathbb{Z}[X \sqcup Y ]/ (I + J) $.

Moreover, given two ring homomorphisms from $ \mathbb{Z}[X] $ and $ \mathbb{Z}[Y] $,
this uniquely defines a ring homomorphism from $ \mathbb{Z}[X \sqcup Y] $ using the universal
property in Problem 2. Moreover, this ring homomorphism factors through the canonical injections
since it acts the same way on the basis. We thus conclude our proof.

\section{Problem 5}

Let $ \mathcal{P} $ be the set of all prime ideals in $ R $.

Order $ \mathcal{P} $ by reverse inclusion. To apply Zorn's Lemma, we should show
that every chain has a maximal (which in this case means minimal using set inclusion)
element.

Let $ P_{1} \supseteq P_{2} \supseteq \ldots $ be chain in $ \mathcal{P} $.
Take the intersection of all the ideals in the chain.
This gives another prime ideal contained in all the ideals in the chain.
Thus, by Zorn's Lemma, we conclude the proof.



\end{document}
